\documentclass[11pt,a4paper]{article}
% Estilo comum dos PDFs individuais da Sprint 1.
\usepackage[utf8]{inputenc}
\usepackage[T1]{fontenc}
\usepackage[brazilian]{babel}
\usepackage[a4paper,margin=2.35cm]{geometry}
\usepackage{xcolor}
\usepackage{hyperref}
\usepackage{enumitem}
\usepackage{listings}
\usepackage{parskip}
\usepackage{needspace}
\definecolor{azul}{RGB}{20,63,112}
\hypersetup{colorlinks=true,linkcolor=azul,urlcolor=azul}
\setlist[itemize]{leftmargin=*,itemsep=3pt}
\lstset{basicstyle=\ttfamily\small,frame=single,breaklines=true,columns=fullflexible}
\sloppy
\newcommand{\codigo}[1]{\texttt{#1}}
\newcommand{\titulo}[2]{\begin{center}{\LARGE\bfseries #1\par}\vspace{4pt}{\large #2\par}\vspace{8pt}\rule{0.8\textwidth}{0.4pt}\par\vspace{5pt}{\small Matheus Moreli Durão Silva -- Matrícula 12621CBS217\par Cibersegurança - Bacharelado - Noturno\par}\end{center}}
\newcommand{\cardfonte}[1]{\par\smallskip\noindent\textit{Referência: }\href{#1}{Trello - enunciado e checklist do card}.}

\begin{document}\titulo{S1\_17 - Versionamento colaborativo}{Branch, merge e Pull Request}
\section*{Entregável correto}
O card pede uma branch criada, integrada e documentada no GitHub. Uma branch é uma linha alternativa de trabalho: ela permite desenvolver uma mudança sem alterar imediatamente a versão principal, chamada \codigo{main}.
\section*{Fluxo explicado}
Criar uma branch com \codigo{git switch -c feature/validacao-json} copia o ponto atual para uma linha de trabalho. Nela, criam-se commits da funcionalidade. Ao terminar, a branch pode ser enviada ao GitHub com \codigo{git push -u origin feature/validacao-json}. A Pull Request (PR) mostra a comparação entre a branch e \codigo{main}; após revisão, o merge integra as alterações.
\begin{lstlisting}
git switch -c feature/validacao-json
git add app.py
git commit -m "feat: validar formato JSON"
git push -u origin feature/validacao-json
\end{lstlisting}
Conflito acontece quando duas versões alteram a mesma região de forma incompatível. Git marca as alternativas, a pessoa decide qual conteúdo fica, salva, adiciona o arquivo resolvido e conclui o merge. Não é um erro ``misterioso''; é uma decisão que precisa ser explícita.
\section*{Evidência no projeto}
O repositório possui a branch \codigo{feature/validacao-json}, posteriormente integrada por Pull Request. A documentação correta da etapa aponta a finalidade da branch, os commits, a PR e o resultado do merge, em vez de apenas dizer que ``deu certo''.
\section*{Checklist demonstrado}
Versionamento, branches, commits, merge, conflitos, Pull Requests e colaboração no GitHub. Para revisar, comparar a PR com a \codigo{main} e confirmar que os arquivos desejados foram integrados.
\cardfonte{https://trello.com/c/Ec5uKKhf}
\end{document}
